%%% FIELD vector-definition-construction
Choose \(m\to\infty\) with \(m/n\to0\). Conditionally on the observed full sample \(O_1,\ldots,O_n\), draw \(B\) according to \(\mathbb P_B\), uniformly over all subsets of \(\{1,\ldots,n\}\) having cardinality \(m\). For every realized \(B\), refit all folds, generalized-inverse scales, order-statistic clips, preliminary quantiles, endpoint corrections, and densities. For \(a\in\mathcal A\) and \(\tau\in\mathcal T\), set \(V_{B,a}(\tau)=m\widehat H_{a,m,B}^{-1}\{\widehat q_{a,B}^{\mathrm{bc}}(\tau)-\widehat q_a^{\mathrm{bc}}(\tau)\}\), set \(T_{B,n}(\tau)=\widehat\ell_{1,n}V_{B,1}(\tau)-\widehat\ell_{0,n}V_{B,0}(\tau)\), and define
\[
c_{n,p}(\tau)=\inf\{t\in\mathbb R:\mathbb P_B\{T_{B,n}(\tau)\le t\mid O_1,\ldots,O_n\}\ge p\}.
\]
For simultaneous inference over the finite set \(\mathcal T\), define
\[
M_{B,n}=\max_{\tau\in\mathcal T}|T_{B,n}(\tau)|,
\qquad
c^{\max}_{n,p}=\inf\{t\in\mathbb R:\mathbb P_B\{M_{B,n}\le t\mid O_1,\ldots,O_n\}\ge p\}.
\]
Thus the conditional law is the complete uniform empirical subsampling measure; conditionally independent draws from \(\mathbb P_B\) give its Monte Carlo approximation.
%%% FIELD anti-concentration-condition
Over the gamma-separated laws \(P\in\mathcal P_{\mathrm S}\) for which the common deterministic heavy- and Gaussian-arm schedules apply at sample size \(n\), the moving row laws satisfy both
\[
\lim_{\eta\downarrow0}\limsup_{n\to\infty}\sup_P
\max_{\tau\in\mathcal T}\sup_{x\in\mathbb R}
P_P\{|Z_{P,n}(\tau)-x|\le\eta\}=0
\]
and
\[
\lim_{\eta\downarrow0}\limsup_{n\to\infty}\sup_P
\sup_{x\in\mathbb R}
P_P\{\bigl|\|Z_{P,n}\|_\infty-x\bigr|\le\eta\}=0.
\]
%%% FIELD vector-theorem-statement
Fix deterministic learner schedules satisfying the heavy- and Gaussian-arm hypotheses of \(\text{thm:joint-qte}\) both at \(N=n\) and after every refit at \(N=m\), including \(J_N\log N/N\to0\), \(J_N^{-s}\log(N/J_N)\to0\), and every displayed approximation-tail condition, where \(m\to\infty\) and \(m/n\to0\). Under \(\text{ass:gamma-separation}\) and \(\text{ass:anti-concentration}\), for every \(\varepsilon>0\),
\[
\sup_P P_P\!\left\{
d_{\mathrm{BL}}\!\left(
\mathcal L_B(T_{B,n}\mid O_1,\ldots,O_n),
\mathcal L_P(Z_{P,n})\right)>\varepsilon\right\}\longrightarrow0,
\]
where the supremum is over exactly the gamma-separated laws \(P\in\mathcal P_{\mathrm S}\) on which those common schedules apply. Moreover, the asymmetric pointwise intervals obey
\[
\max_{\tau\in\mathcal T}\sup_P\left|
P_P\!\left\{\Delta_P(\tau)\in
\left[
\widehat\Delta^{\mathrm{bc}}(\tau)-\frac{\widehat H_n}{n}c_{n,1-\alpha/2}(\tau),
\widehat\Delta^{\mathrm{bc}}(\tau)-\frac{\widehat H_n}{n}c_{n,\alpha/2}(\tau)
\right]\right\}-(1-\alpha)\right|\longrightarrow0.
\]
Define the finite-\(\mathcal T\) simultaneous band by
\[
\mathcal C_n^{\mathrm{sim}}(\tau)=
\left[
\widehat\Delta^{\mathrm{bc}}(\tau)-\frac{\widehat H_n}{n}c^{\max}_{n,1-\alpha},
\widehat\Delta^{\mathrm{bc}}(\tau)+\frac{\widehat H_n}{n}c^{\max}_{n,1-\alpha}
\right].
\]
Then
\[
\sup_P\left|
P_P\!\left\{\Delta_P(\tau)\in\mathcal C_n^{\mathrm{sim}}(\tau)
\text{ for every }\tau\in\mathcal T\right\}-(1-\alpha)
\right|\longrightarrow0.
\]
For every fixed noncritical \(P\in\mathcal P_{\mathrm S}\) satisfying the same schedules, the conditional law converges in \(P\)-probability along any subsequence with \(\ell_{a,P,n}\to\ell_{a,P}\) to the law of \(\ell_{1,P}L_{1,P}-\ell_{0,P}L_{0,P}\), and the pointwise intervals and simultaneous band have limiting coverage \(1-\alpha\). There exist gamma-separated two-heavy row sequences with \(\gamma_{0,P_n}-\gamma_{1,P_n}=c/\log n\) for which raw scalar subsampling at \(m=n^\beta\), \(0<\beta<1\), converges to a different arm mixture from the full-sample moving row law. This counterexample refutes uniform approximation by a single fixed-law limit and raw scalar subsampling, but not the moving-law validity of the refitted arm-vector procedure.
%%% FIELD vector-theorem-proof
We prove conditional approximation to the moving law in \(\text{def:moving-row-law}\), uniformly over the scheduled gamma-separated domain. At sample size \(N\), put
\[
W_{P,N}^{\circ}=\left\{
\frac N{\widehat H_{a,N}}(\widehat q_{a,N}^{\mathrm{bc}}-q_{a,P})
\right\}_{a\in\mathcal A}.
\tag{S1}
\]
The uniform armwise expansions and scale consistency in \(\text{thm:oracle-preservation}\) and \(\text{thm:gaussian-branch}\), together with the gamma separation in \(\text{ass:gamma-separation}\), yield a permutation-symmetric leading statistic \(S_{P,N}\) such that, for every \(\varepsilon>0\),
\[
\sup_PP_P\{\|W_{P,N}^{\circ}-S_{P,N}\|>\varepsilon\}\to0,
\qquad
\sup_Pd_{\mathrm{BL}}\{\mathcal L_P(S_{P,N}),\mathcal L_P(L_P)\}\to0.
\tag{S2}
\]
Here and below every supremum is over the gamma-separated laws in \(\mathcal P_{\mathrm S}\) on which the displayed schedules apply. For \(\gamma_a=\infty\), \(\text{ass:boundary-link}\) and the common Hölder radius in \(\text{ass:holder}\) imply a common positive lower bound for \(\pi_a\) on the compact chart. Hence the upper \(k_N/N\) quantile of \(1/\widehat\pi_a\) lies uniformly within \(o_P(1)\) of its finite endpoint, its Hill log-spacing average is \(o_P(1)\), and \(\widehat\gamma_a\to_P\infty\). Thus the classification used in (S2) also holds on the bounded-support branch.

Let \(g:\mathbb R^{2R}\to[-1,1]\) be one-Lipschitz and define the complete-subsample average
\[
U_{n,m,P}(g)=\binom nm^{-1}\sum_{|b|=m}g(S_{P,b})
=E_B\{g(S_{P,B})\mid O_1,\ldots,O_n\}.
\tag{S3}
\]
Successively subtracting conditional expectations over proper coordinate subsets gives the orthogonal decomposition
\[
g(S_{P,m})-E_Pg(S_{P,m})
=\sum_{j=1}^m\sum_{|C|=j}h_{P,m,j}(O_C),
\tag{S4}
\]
where each \(h_{P,m,j}\) has conditional mean zero given any \(j-1\) arguments. If
\(\zeta_{P,m,j}=E_Ph_{P,m,j}(O_1,\ldots,O_j)^2\), independence in \(\text{ass:iid}\), orthogonality, and counting the \(m\)-subsets containing a fixed \(j\)-subset give
\[
\begin{split}
\operatorname{Var}_P\{U_{n,m,P}(g)\}
&=\sum_{j=1}^m\frac{\binom mj^2}{\binom nj}\zeta_{P,m,j}\\
&\le\frac mn\sum_{j=1}^m\binom mj\zeta_{P,m,j}
=\frac mn\operatorname{Var}_P\{g(S_{P,m})\}
\le4m/n.
\end{split}
\tag{S5}
\]
The inequality uses \(\binom mj/\binom nj\le m/n\), while the last bound uses \(|g|\le1\). Since a uniform \(m\)-subset of an iid sample has unconditionally the law of \(m\) iid observations, (S2)--(S5) imply
\[
U_{n,m,P}(g)-E_Pg(L_P)\to_P0
\tag{S6}
\]
uniformly in \(P\).

The family \(\{L_P\}\) is uniformly tight. Indeed, the Gaussian coordinates have a common second-moment bound by \(\text{ass:finite-variance}\) and the uniform influence-function expansion in \(\text{thm:gaussian-branch}\), while the heavy coordinates have the uniform small- and large-jump bounds established in \(\text{thm:oracle-preservation}\). On a common compact ball, the unit bounded-Lipschitz class has a finite uniform supremum-norm net. Applying (S5) to this net and using uniform tightness off the ball upgrades (S6) to
\[
\sup_PP_P\!\left\{
d_{\mathrm{BL}}\bigl(\mathcal L_B(S_{P,B}\mid O_1^n),
\mathcal L_P(L_P)\bigr)>\varepsilon\right\}\longrightarrow0.
\tag{S7}
\]
The unconditional-law identity and (S2) also give
\[
\begin{split}
&E_P\mathbb P_B\{\|W_{P,B}^{\circ}-S_{P,B}\|>\varepsilon\mid O_1^n\}\\
&\qquad=P_P\{\|W_{P,m}^{\circ}-S_{P,m}\|>\varepsilon\}\longrightarrow0
\end{split}
\tag{S8}
\]
uniformly. Markov's inequality applied to the conditional error probability in (S8) combines (S7)--(S8) into conditional uniform convergence of the refitted arm vector.

By \(\text{def:vector-subsampling}\),
\[
V_{B,a}=W_{P,B,a}^{\circ}-C_{B,a},
\qquad
C_{B,a}=\frac m{\widehat H_{a,m,B}}
(\widehat q_{a,n}^{\mathrm{bc}}-q_{a,P}).
\tag{S9}
\]
Uniform tightness in (S2), scale consistency, and Hall inversion under \(\text{ass:hall}\) give
\[
\|C_{B,a}\|=O_P\!\left(\frac m{H_{a,P,m}}\frac{H_{a,P,n}}n\right)
=\begin{cases}
O_P\{(m/n)^{1-1/\gamma_{a,P}}\},&a\in\mathcal A_{\mathrm H},\\
O_P\{(m/n)^{1/2}\},&a\in\mathcal A_{\mathrm G}.
\end{cases}
\tag{S10}
\]
Both rates are uniformly \(o_P(1)\): the Gaussian conclusion uses \(m/n\to0\), while in the heavy branch \(\text{ass:gamma-separation}\) gives
\(1-1/\gamma_{a,P}\ge\epsilon_\gamma/(1+\epsilon_\gamma)>0\). Scale consistency also gives
\(\max_a|\widehat\ell_{a,n}-\ell_{a,P,n}|=o_P(1)\) uniformly. Recombination has Lipschitz norm at most two because \(0\le\widehat\ell_{a,n},\ell_{a,P,n}\le1\). Therefore (S7)--(S10) prove
\[
\sup_PP_P\!\left\{
d_{\mathrm{BL}}\!\left(
\mathcal L_B(T_{B,n}\mid O_1^n),
\mathcal L_P(Z_{P,n})\right)>\varepsilon\right\}\longrightarrow0.
\tag{S11}
\]

Let
\[
X_{P,n}=\frac n{\widehat H_n}(\widehat\Delta^{\mathrm{bc}}-\Delta_P).
\tag{S12}
\]
The full-sample approximation in \(\text{thm:joint-qte}\), uniform consistency of \(\widehat H_n/H_{P,n}\), and uniform tightness imply
\[
\sup_Pd_{\mathrm{BL}}\{\mathcal L_P(X_{P,n}),\mathcal L_P(Z_{P,n})\}\to0.
\tag{S13}
\]
For a real random variable \(U\), approximate the half-line indicator \(\mathbf1\{U\le x\}\) from above and below by piecewise-linear functions with transition width \(\eta\). After rescaling, these functions have bounded-Lipschitz norm at most \(1+1/\eta\). Consequently a bounded-Lipschitz discrepancy \(\delta\) changes the half-line probability by at most \((1+1/\eta)\delta\) plus the probability assigned by the reference law to an interval of radius \(\eta\). Apply this inequality to each coordinate of (S11) and (S13). The coordinatewise modulus in \(\text{ass:anti-concentration}\), followed first by \(n\to\infty\) and then by \(\eta\downarrow0\), yields
\[
\max_{\tau\in\mathcal T}\sup_P\left|
P_P\!\left\{
c_{n,\alpha/2}(\tau)\le X_{P,n}(\tau)
\le c_{n,1-\alpha/2}(\tau)\right\}-(1-\alpha)
\right|\to0.
\tag{S14}
\]
Convergence in probability of the conditional bounded-Lipschitz discrepancy suffices here because that discrepancy is bounded, so truncation converts it to uniform convergence of its expectation. The event in (S14) is exactly coverage by the asymmetric pointwise interval in the statement.

For simultaneous coverage, use the map \(\phi:\mathbb R^R\to\mathbb R\), \(\phi(x)=\|x\|_\infty\). It is one-Lipschitz because
\[
|\|x\|_\infty-\|y\|_\infty|\le\|x-y\|_\infty.
\tag{S15}
\]
Hence composition with \(\phi\) cannot increase bounded-Lipschitz distance. Equations (S11) and (S13) therefore compare, uniformly, the conditional law of
\(M_{B,n}=\phi(T_{B,n})\), the law of \(\phi(X_{P,n})\), and the row law of \(\phi(Z_{P,n})\). Applying the same two-sided half-line smoothing inequality, now with the max-law modulus in \(\text{ass:anti-concentration}\), gives
\[
\sup_P\left|
P_P\{\|X_{P,n}\|_\infty\le c^{\max}_{n,1-\alpha}\}-(1-\alpha)
\right|\to0.
\tag{S16}
\]
By the definitions of \(X_{P,n}\) and \(\mathcal C_n^{\mathrm{sim}}\), the event in (S16) equals
\[
\{\Delta_P(\tau)\in\mathcal C_n^{\mathrm{sim}}(\tau)
\text{ for every }\tau\in\mathcal T\}.
\tag{S17}
\]
This proves the simultaneous coverage assertion. The finiteness of \(\mathcal T\) ensures that the statistic is a finite-dimensional continuous maximum and that the displayed band is a finite collection of intervals.

For fixed noncritical \(P\), every subsequence has a further subsequence along which both bounded recombination weights converge. Equation (S11) then gives the stated fixed-law conditional limit; (S14) and (S16) give the fixed-law pointwise and simultaneous coverage conclusions. To certify the scope of the fixed-limit failure, take the two-heavy family in \(\text{prop:smooth-inclusion}\), an interior \(\gamma\in(1,2)\), equal Hall constants, fixed nondegenerate endpoint score laws and target densities, and
\[
\gamma_{0,P_n}=\gamma+\frac c{2\log n},
\qquad
\gamma_{1,P_n}=\gamma-\frac c{2\log n},
\qquad c>0.
\tag{S18}
\]
These rows remain gamma-separated with common constants. Hall inversion yields
\[
\log\{h_{0,P_n,n}/h_{1,P_n,n}\}
=\{\gamma_{0,P_n}^{-1}-\gamma_{1,P_n}^{-1}\}\log n+o(1)
\longrightarrow-c/\gamma^2.
\tag{S19}
\]
Thus the full-sample row law retains the mixture
\(L_{1,P_n}-e^{-c/\gamma^2}L_{0,P_n}\), while fixing a row and sending a new sample size to infinity eliminates the larger-index arm. Independence and nondegeneracy distinguish these laws: equality of their characteristic functions near zero would force the characteristic function of a nonzero multiple of \(L_{0,P_n}\) to be one near zero, hence force degeneracy.

If raw scalar subsampling uses \(m=n^\beta\), \(0<\beta<1\), then
\[
\log\{h_{0,P_n,m}/h_{1,P_n,m}\}\longrightarrow-\beta c/\gamma^2,
\tag{S20}
\]
so it approaches \(L_1-e^{-\beta c/\gamma^2}L_0\), not the full-sample moving mixture. The refitted vector instead standardizes each arm at its own \(m\)-scale and recombines with the full-sample \(\widehat\ell_{a,n}\), reproducing (S19). Hence (S18)--(S20) refute only a common fixed-limit approximation and raw scalar subsampling; they affirm the moving-law validity in (S11).
